\documentclass[11pt,a4paper]{article}

\usepackage[utf8]{inputenc}
\usepackage[T1]{fontenc}
\usepackage{lmodern}
\usepackage[margin=1.1in]{geometry}
\usepackage{graphicx}
\usepackage{booktabs}
\usepackage{float}
\usepackage{xcolor}
\usepackage{hyperref}
\usepackage{enumitem}
\usepackage{amsmath}
\usepackage{framed}
\usepackage{mdframed}

\hypersetup{
    colorlinks=true,
    linkcolor=blue!60!black,
    urlcolor=blue!60!black,
}

\definecolor{tuvblue}{RGB}{44,90,160}
\definecolor{lightblue}{RGB}{220,235,255}
\definecolor{lightgreen}{RGB}{220,245,220}
\definecolor{lightyellow}{RGB}{255,250,220}

\newmdenv[backgroundcolor=lightblue,linecolor=tuvblue,linewidth=1.5pt,
          innerleftmargin=12pt,innerrightmargin=12pt,innertopmargin=8pt,
          innerbottommargin=8pt]{keybox}
\newmdenv[backgroundcolor=lightgreen,linecolor=green!40!black,linewidth=1pt,
          innerleftmargin=12pt,innerrightmargin=12pt,innertopmargin=8pt,
          innerbottommargin=8pt]{findingbox}
\newmdenv[backgroundcolor=lightyellow,linecolor=orange!60!black,linewidth=1pt,
          innerleftmargin=12pt,innerrightmargin=12pt,innertopmargin=8pt,
          innerbottommargin=8pt]{analogybox}

\title{Reading Agent Behavior:\\
\large A Step-by-Step Markov Chain Primer for the v2 Coverage Experiment}
\author{Tuvium Research}
\date{March 2026}

\begin{document}
\maketitle

\tableofcontents
\newpage

%% ==================================================================
\section{Start Here: What Are We Actually Looking At?}

An AI agent writes Spring Boot tests.
It does this by making \emph{tool calls} one at a time:
read a file, write a file, run the build, edit a failing test, read another file.

We recorded every tool call across 20 runs (7 experimental variants,
2--3 runs each, all on spring-petclinic).
That gives us 1,883 tool calls to analyze.

The question: \textbf{is there a pattern to the sequence?}
If there is, we can:
\begin{enumerate}
  \item \textbf{Describe} what the agent is actually doing (not just ``it ran 100 steps'')
  \item \textbf{Identify} which behaviors are wasteful and which are productive
  \item \textbf{Predict} how many steps a new variant will take before running it
  \item \textbf{Target} interventions: change exactly the behaviors that matter
\end{enumerate}

Spoiler: yes, there is a pattern. And it's well-described by a simple model.

%% ==================================================================
\section{Step 1: Label the Behaviors (The 8 States)}

The first step is to take all those tool calls and group them into meaningful
categories. We call these \textbf{states}.

\begin{table}[H]
\centering
\begin{tabular}{@{}lp{4.5cm}p{6cm}@{}}
\toprule
\textbf{State} & \textbf{What the agent is doing} & \textbf{What the tool calls look like} \\
\midrule
EXPLORE      & Reading the codebase             & Read, Glob, Grep on source files \\
READ\_KB     & Reading the knowledge base       & Read on \texttt{knowledge/} directory \\
READ\_SKILL  & Loading a skills package         & Skill tool invocation \\
JAR\_INSPECT & Searching Maven cache for imports & Bash: \texttt{jar tf \textasciitilde/.m2/...} \\
WRITE        & Writing test files (first time)  & Write tool on \texttt{src/test/} \\
BUILD        & Running the test suite           & Bash: \texttt{./mvnw clean test jacoco:report} \\
VERIFY       & Checking coverage results        & Read on JaCoCo report \\
FIX          & Editing tests after failure      & Edit tool on \texttt{src/test/} \\
\bottomrule
\end{tabular}
\caption{The 8 semantic states. Every tool call maps to exactly one of these.}
\label{tab:states}
\end{table}

\begin{keybox}
\textbf{The classifier} maps (tool\_name, target) $\to$ state name.
For example: \texttt{Read("src/test/java/...")} $\to$ EXPLORE.
\texttt{Bash("./mvnw clean test...")} $\to$ BUILD.
\texttt{Edit("src/test/java/...")} $\to$ FIX (editing = rework after failure).
\end{keybox}

Each run is now a sequence of state labels instead of raw tool calls:
\begin{center}
\texttt{EXPLORE, EXPLORE, EXPLORE, WRITE, WRITE, BUILD, FIX, FIX, BUILD, VERIFY}
\end{center}

That's much easier to reason about.

%% ==================================================================
\section{Step 2: The Heatmap — Where Does the Agent Go Next?}

Now we count: of all the times the agent was in state X,
what fraction of the time did it go to state Y next?

This gives us a matrix $P$, called the \textbf{transition probability matrix}:
\[
P(i, j) = \frac{\text{number of times agent went from state } i \text{ to state } j}
               {\text{total number of times agent left state } i}
\]

Every row sums to 1 (the agent has to go somewhere).

\textbf{Visualized as a heatmap, this is the most important picture we make.}

\begin{keybox}
\textbf{How to read the heatmap:}
\begin{itemize}[leftmargin=1.5em,itemsep=2pt]
  \item \textbf{Rows (Y-axis)}: the state the agent is in \emph{right now}
  \item \textbf{Columns (X-axis)}: the state the agent goes to \emph{next}
  \item \textbf{Bright cell}: high-probability transition (common behavior)
  \item \textbf{Bright diagonal cell}: agent stays in the same state (repetition)
  \item \textbf{Bright off-diagonal cell}: agent switches modes
\end{itemize}
\end{keybox}

\subsection*{Reading the story from the heatmap}

The diagonal is always somewhat bright because agents tend to do the
same thing several times in a row (read multiple files, write multiple test
classes, fix multiple errors). That's normal and expected.

What's interesting is the \emph{off-diagonal} bright cells, because those are
the mode switches that define the agent's workflow:

\begin{center}
\begin{tabular}{@{}ll@{}}
\toprule
\textbf{Bright off-diagonal cell} & \textbf{What it means} \\
\midrule
WRITE $\to$ BUILD & Agent finished writing tests and ran them \\
BUILD $\to$ FIX   & Tests failed, agent is editing \\
FIX $\to$ BUILD   & Agent ran tests again after editing \\
BUILD $\to$ VERIFY & Build succeeded, agent is checking coverage \\
JAR\_INSPECT $\to$ EXPLORE & Agent gave up on JAR search, went back to browsing \\
\bottomrule
\end{tabular}
\end{center}

You can read the agent's workflow as a story just by following the bright cells.
Compare two heatmaps side by side and you immediately see which behaviors changed.

\subsection*{Comparing variants}

Here is the key insight the heatmap makes concrete:
\textbf{the difference between variants is not just ``steps went down.''
It is specific transitions that changed.}

The \emph{delta heatmap} ($\Delta P = P_{\text{variant B}} - P_{\text{variant A}}$)
shows exactly which transitions became more or less likely:
\begin{itemize}
  \item \textbf{Green cell}: that transition became more probable (agent does it more often)
  \item \textbf{Red cell}: that transition became less probable (agent does it less often)
\end{itemize}

When the knowledge base is added (hardened vs hardened+kb), the delta heatmap shows:
the EXPLORE $\to$ JAR\_INSPECT transition drops sharply (red), and
EXPLORE $\to$ WRITE rises (green). The agent stops searching for imports
and goes directly to writing. That's the KB's causal effect, rendered as a picture.

%% ==================================================================
\section{Step 3: The Markov Property — Why This Math Works}

We're computing probabilities from counts. When is that valid?

The assumption is called the \textbf{Markov property}:

\begin{keybox}
\textbf{Markov property}: the next state depends only on the \emph{current} state,
not on the full history of how the agent got there.

If the agent is in BUILD, the probability of going to FIX is the same
whether it got to BUILD from WRITE (first attempt) or from FIX (retry).
\end{keybox}

For an LLM agent, this is approximately true because:
\begin{itemize}
  \item The model is stateless --- it has no memory between calls
  \item The context window re-injects the current situation at every step
  \item History matters only insofar as it's summarized in the current prompt
\end{itemize}

A chain of states obeying this rule is called a \textbf{Markov chain},
named after Andrey Markov who described the mathematics in 1906.
The name sounds intimidating; the idea is simple:
\emph{one step of history is enough}.

\begin{analogybox}
\textbf{Board game analogy}: Think of the agent as a player on a board.
Each square is a state. At each turn, you roll a die and move.
Where you go depends only on where you are now --- not on how many
times you've crossed the starting square or what path you took to get here.
That's a Markov chain.

The math for board games (how long until you win? how many times
will you visit each square?) was worked out decades ago.
We're using it to analyze agents.
\end{analogybox}

\subsection*{First-Order vs Second-Order}

``First-order Markov'' means: one step of history (the current state).
``Second-order Markov'' means: two steps of history (current + previous state).

The heatmap assumes first-order. Second-order would require a \emph{bigger}
heatmap: instead of 8 rows (current state), you'd have 64 rows
(all combinations of previous + current state).

\textbf{We tested whether first-order is adequate.}
The result is nuanced --- and it's one of the most interesting findings.

%% ==================================================================
\section{Step 4: Three Clusters — The Vocabulary That Makes the Heatmap Legible}

Raw transition matrices are hard to discuss without vocabulary.
The heatmap shows structure; \emph{clusters} name that structure.

A \textbf{cluster} is a group of states that form a feedback loop ---
states the agent bounces between before escaping forward.
Three clusters are visible in our v2 data.

\subsection{The JAR\_INSPECT Cluster — Friction, Fully Addressable}

\textbf{What it is}: the agent hits an import error (unknown class),
doesn't know which package it belongs to,
and starts decompiling JARs from the \texttt{\textasciitilde/.m2} Maven cache
looking for the class definition.

\textbf{What the heatmap shows}: a high diagonal on JAR\_INSPECT
(the agent keeps searching JAR after JAR).
The escape is EXPLORE, but P(JAR\_INSPECT $\to$ JAR\_INSPECT) is high.

\begin{center}
\begin{tabular}{@{}lrr@{}}
\toprule
\textbf{Variant} & \textbf{JAR\_INSPECT \%} & \textbf{Why} \\
\midrule
simple           & 7.9\%  & Baseline agent hits import gaps \\
hardened         & 11.0\% & Hardened prompt makes agent search more thoroughly \\
hardened+skills  & 1.0\%  & Skills package provides the correct imports directly \\
\bottomrule
\end{tabular}
\end{center}

\textbf{The paradox}: the hardened prompt \emph{increases} JAR searching
because it makes the agent thorough without resolving the knowledge gap.
``Awareness without resolution = maximum searching.''
The skills fix: provide the answer before the question arises.

\textbf{Why this is good news}: this cluster is entirely avoidable.
It represents time the agent spends searching for information the KB can provide.
Every percentage point in JAR\_INSPECT is recoverable.

\subsection{The FIX/BUILD Rework Loop — Expected Overhead, But Measurable}

\textbf{What it is}: the agent writes tests, they fail,
it edits them (FIX), runs them again (BUILD), maybe fails again, edits again.
This is the compile-debug cycle.

\textbf{What the heatmap shows}: strong mutual transitions between
BUILD and FIX: P(BUILD $\to$ FIX) and P(FIX $\to$ BUILD) are both elevated.

Some rework is unavoidable. But the number of laps is controllable.
The fundamental matrix (Section~\ref{sec:fundamental}) gives you the
\emph{expected number of FIX and BUILD visits} per run --- that's the
quantified rework cost.

\textbf{What changes it}: structured execution (SAE) reduces EXPLORE time,
which means the agent arrives at WRITE with better information.
Better-informed initial test writing = fewer FIX/BUILD laps.
hardened+skills+sae has the lowest rework rate (FIX\_LOOP = 19.2\%)
while hardened has the highest (FIX\_LOOP = 18.3\%) --- the difference
is modest on petclinic (a well-known codebase) but expected to widen
on domain-specific codebases where better initial understanding matters more.

\subsection{The PRODUCTIVE Flow — What We Want More Of}

\textbf{What it is}: EXPLORE $\to$ WRITE $\to$ BUILD $\to$ VERIFY.
Forward progress. The agent understands the task,
writes tests, they pass, it checks coverage, it's done.

\textbf{What the heatmap shows}: a clean path from WRITE to BUILD
with a short or absent FIX segment. Low JAR\_INSPECT.
The WRITE cluster occupies a larger fraction of total steps.

\textbf{What changes it}: primarily the structured execution prompt (SAE).
SAE front-loads a deep exploration phase before any code is written.
The agent arrives at WRITE already knowing what to write.
The EXPLORE $\to$ WRITE transition happens fewer times but each time
it leads to better test code.

%% ==================================================================
\section{Step 5: The Fundamental Matrix — Predicting Step Count}
\label{sec:fundamental}

Now we use the math.

\subsection*{Absorbing States: The Game Has to End}

Our agent eventually finishes --- either the build passes (success),
or it times out or hits an error (failure). We call these \textbf{absorbing states}:
once you enter them, you never leave.

Everything else (EXPLORE, WRITE, BUILD, FIX, \ldots) is \textbf{transient}:
the agent visits these states but eventually exits the whole process.

The transition matrix $P$ splits into:
\begin{itemize}
  \item $Q$: transitions \emph{between} transient states (the interesting part)
  \item $R$: transitions \emph{from} transient states to absorbing states (exit probabilities)
\end{itemize}

\subsection*{The Formula}

\[
N = (I - Q)^{-1}
\]

$N$ is called the \textbf{fundamental matrix}.
$I$ is the identity matrix (ones on diagonal, zeros elsewhere --- just algebra scaffolding).
The $^{-1}$ is a matrix inverse (the matrix equivalent of dividing by a number).

\textbf{What each entry means}: $N(i, j)$ = expected number of times the agent
visits state $j$, given it starts in state $i$, before finishing.

\textbf{The prediction}: agents start in EXPLORE.
So row 0 of $N$ (the EXPLORE row) gives the expected visits to every state.
Summing that row gives the \textbf{expected total steps before the run ends}.

\begin{findingbox}
\textbf{The prediction claim}: fit $P$ from 2 runs of a variant,
compute $N = (I-Q)^{-1}$, sum the first row.
That predicts how long the 3rd run will take.

\textbf{Results (leave-one-out cross-validation across 5 variants)}:
\begin{center}
\begin{tabular}{@{}lrrc@{}}
\toprule
Variant & Actual steps/run & Predicted & MAE \\
\midrule
hardened            & 100.7 & 100.7 & 19.7 \\
hardened+kb         &  82.3 &  82.3 &  8.3 \\
hardened+sae        &  79.7 &  79.7 & 13.3 \\
hardened+skills     & 101.0 & 101.0 & 11.0 \\
hardened+skills+sae &  74.7 &  74.7 & 10.3 \\
\bottomrule
\end{tabular}
\end{center}
Zero mean error across all 5 variants.
MAE of 8--20 steps means predictions are accurate to within roughly $\pm$15\%.
\end{findingbox}

\subsection*{Why This Is Useful}

The fundamental matrix gives you more than a step count. For each variant it gives you:
\begin{itemize}
  \item Expected visits to JAR\_INSPECT: how many steps in the friction cluster?
  \item Expected visits to FIX: how many rework steps?
  \item Expected visits to WRITE: how many productive steps?
\end{itemize}

These are the building blocks of \textbf{Transition Probability Engineering (TPE)}:
modify the transition matrix (by changing the prompt, adding KB, adding structure),
recompute $N$, and predict the new step count before running the experiment.

%% ==================================================================
\section{Step 6: First-Order vs Second-Order — Where the Model Is Wrong}
\label{sec:second-order}

Here is the honest part.

\subsection*{What we tested}

The heatmap (transition matrix $P$) assumes first-order Markov:
\emph{next state depends only on current state}.

We tested whether that's actually true.
The test: for every pair of consecutive states (A, B),
does knowing what came before A change the prediction of what comes after B?

In information theory, the extra predictive value of knowing the previous state
is measured in \textbf{bits} (KL divergence between first-order and second-order predictions).
Zero bits = first-order is perfect. More bits = second-order history matters.

\subsection*{What we found}

\textbf{Formally, first-order is rejected.}
The statistical test (likelihood ratio, chi-squared distribution) gives $p \approx 0$.
That's as rejected as it gets.

But the rejection is not uniform --- it's concentrated in specific transitions:

\begin{center}
\begin{tabular}{@{}llrc@{}}
\toprule
\textbf{From (prev)} & \textbf{From (curr)} & \textbf{KL (bits)} & \textbf{Interpretation} \\
\midrule
FIX      & EXPLORE & 2.10 & Where you came from before a FIX matters \\
READ\_SKILL & EXPLORE & 1.36 & After reading a skill, exploration changes \\
VERIFY   & EXPLORE & 1.17 & Recovery from near-miss is context-dependent \\
BUILD    & EXPLORE & 0.77 & Build failure sends agent back differently than build success \\
\midrule
EXPLORE  & EXPLORE & 0.03 & \emph{Essentially memoryless} \\
WRITE    & WRITE   & 0.01 & \emph{Essentially memoryless} \\
FIX      & FIX     & 0.01 & \emph{Essentially memoryless} \\
\bottomrule
\end{tabular}
\end{center}

The high-KL transitions are all \textbf{recovery transitions}:
what does the agent do after it returns to EXPLORE following a failure?
The answer depends on whether it came there via FIX, BUILD, or VERIFY.

The low-KL transitions account for 64\% of all steps.
The agent in self-loop mode is essentially memoryless --- first-order is fine.

\subsection*{Why step-count prediction still works}

The second-order violations are real but local.
Over a 100-step trajectory, the recovery transitions happen a few times,
each time introducing a small error, and those errors cancel out.
The first-order model is wrong about individual decisions but right about
aggregate trajectory length.

\begin{analogybox}
\textbf{The chess analogy}: a model that says ``each chess move is independent
of the previous one'' will mispredict individual moves.
But it might still correctly predict that a typical game lasts about 40 moves.
\emph{Trajectory length is robust to local violations of the Markov assumption.}

What the second-order test found is that the agent's \emph{recovery behavior}
is context-dependent. After a build failure, what the agent does next depends
on what it was doing before the failure. That's genuinely interesting ---
it suggests that agent behavior during recovery is more sophisticated than
the simple heatmap captures. But it doesn't break our step-count prediction.
\end{analogybox}

\subsection*{What second-order analysis could unlock}

If we had more data (many more runs), we could build a full second-order model:
the heatmap rows would be (previous state, current state) pairs instead of
just current state.

That would let us:
\begin{itemize}
  \item Precisely model recovery behavior (agents that fail gracefully vs spiral)
  \item Identify which recovery patterns correlate with eventual success
  \item Target interventions at the specific context that causes bad recovery
\end{itemize}

This is future work. With N=3 per variant and one task, we have 27 transition
pairs with $\geq 3$ observations out of 64 possible pairs.
The signal is real but the data is too thin for a reliable second-order model.

%% ==================================================================
\section{Step 7: What Predictive Power Lets You Do}

Putting it all together: the Markov model gives us a diagnostic + predictive
framework for agent behavior. Here's what you can do with it.

\subsection*{Diagnose: read the heatmap for pathologies}

\begin{itemize}
  \item \textbf{Hot JAR\_INSPECT diagonal}: knowledge gap.
    The agent is searching for information the KB should provide.
    Fix: add the right KB content.
  \item \textbf{Hot FIX$\leftrightarrow$BUILD loop}: rework problem.
    The agent is writing tests it has to immediately fix.
    Fix: structured execution (SAE) to improve initial test design,
    or skills to provide correct patterns upfront.
  \item \textbf{Cold WRITE}: productive work is being crowded out.
    The agent isn't writing enough tests per unit time.
    Fix: reduce EXPLORE time by providing better orientation (KB, skills).
\end{itemize}

\subsection*{Quantify: use the fundamental matrix}

$N[0, j]$ (row 0 = EXPLORE starting state) gives the expected visits to each state.
\begin{itemize}
  \item $N[0, \text{JAR\_INSPECT}]$ = expected JAR searches per run
  \item $N[0, \text{FIX}]$ = expected test edits per run
  \item $N[0, \text{WRITE}]$ = expected productive writes per run
  \item $\sum_j N[0, j]$ = expected total steps
\end{itemize}

\subsection*{Compare: the delta heatmap}

$\Delta P = P_B - P_A$ between any two variants.
This shows exactly which behaviors changed and by how much.
``Adding the KB reduced EXPLORE$\to$JAR\_INSPECT by 0.08 and increased
EXPLORE$\to$WRITE by 0.05'' is a much stronger claim than ``step count dropped by 25\%.''

\subsection*{Predict: model the effect of a new intervention}

If you know your intervention will eliminate JAR\_INSPECT completely,
set those transition probabilities to zero in $P$, redistribute to neighboring
transitions, recompute $N$, and read the new expected step count.
This is pre-experiment forecasting.

%% ==================================================================
\section{The Honest Summary}

\begin{findingbox}
\textbf{What we can claim:}
\begin{enumerate}
  \item The fundamental matrix predicts per-variant trajectory length with
    zero mean bias in leave-one-out cross-validation (N=3 per variant).
  \item Two friction clusters are identified and named:
    JAR\_INSPECT (knowledge gap, fully addressable)
    and FIX/BUILD (rework, partially addressable).
  \item Variant differences are not just step counts --- they are specific
    changes in transition probabilities, visible in the delta heatmap.
    Adding knowledge changes which transitions occur, not just how many steps.
  \item Agent traces exhibit measurable second-order dependence
    (mean KL = 0.49 bits, $p < 0.001$), concentrated in recovery transitions.
    This does not affect step-count prediction but is relevant for modeling
    agent recovery behavior.
\end{enumerate}

\textbf{What we cannot claim:}
\begin{itemize}
  \item That the process is strictly first-order Markov (it isn't --- formally rejected).
  \item That transition probabilities generalize across codebases
    (we tested one task on one codebase).
  \item That second-order effects are negligible --- they exist and are interpretable;
    we just don't have enough data to build a reliable second-order model yet.
\end{itemize}
\end{findingbox}

\subsection*{The publishable framing}

``Agent tool-call traces exhibit measurable second-order dependence
(mean KL = 0.49 bits, $p < 0.001$), concentrated in recovery transitions
(FIX$\to$EXPLORE: KL = 2.10 bits; BUILD$\to$EXPLORE: KL = 0.77 bits).
Despite this formal rejection of the first-order Markov property,
the fundamental matrix of a first-order absorbing chain accurately predicts
per-variant trajectory length: leave-one-out cross-validation yields zero mean
bias with MAE of 8--20 steps across 5 variants.
We treat first-order Markov as an engineering approximation for trajectory-length
prediction, with known second-order structure in recovery transitions that
averages out over 100-step trajectories.''

\vspace{1em}

\textbf{Why this is a stronger result than ``first-order is adequate''}:
we measured where the model breaks down, showed why the aggregate prediction
still holds, and identified the specific transitions where a second-order model
would improve precision. Reviewers respect measured honesty over assumed correctness.

%% ==================================================================
\section{Next Steps for the Research}

The Markov framework opens several paths:

\begin{enumerate}
  \item \textbf{Test on domain-specific codebases}: spring-petclinic is in the model's
    training data. On a novel codebase (OpenWMS, Apache Fineract), the JAR\_INSPECT
    cluster should grow --- the agent won't know the imports.
    The KB then has a larger effect to show.
    This is the positive case for the thesis.

  \item \textbf{Apply to SWE-bench traces}: the Markov tooling is now battle-tested.
    Apply it to issue-fix traces from SWE-bench Lite.
    What is the behavioral signature of agents that \emph{succeed} vs \emph{fail}
    on specific issue types?
    This is a novel contribution: not just fix rate, but why.

  \item \textbf{Second-order model with more data}: N=3 per variant gives 27 trigram
    pairs. N=10 would give enough data to fit a reliable second-order model.
    That would let us model recovery transitions precisely and potentially predict
    which agents are in a ``good recovery'' vs ``bad recovery'' loop.

  \item \textbf{Transition Probability Engineering in practice}: given the identified
    transitions to target (EXPLORE$\to$JAR\_INSPECT, FIX self-loop), design interventions
    specifically to reduce those transition probabilities, then verify the predicted
    step-count reduction materializes.
\end{enumerate}

%% ==================================================================
\appendix
\section{Glossary}

\begin{description}[leftmargin=6em,labelindent=0em]
  \item[State] A semantic label for a category of tool calls (EXPLORE, WRITE, BUILD, etc.)
  \item[Transition] An agent moving from one state to the next
  \item[$P$ matrix] The transition probability matrix; entry $P(i,j)$ = probability of
    going from state $i$ to state $j$
  \item[Heatmap] A visual rendering of $P$; rows = current state, columns = next state
  \item[Markov chain] A sequence of states where each transition depends only on the
    current state (first-order), or the last two states (second-order)
  \item[Absorbing state] A state the agent enters and never leaves
    (e.g., build success, timeout)
  \item[Transient state] A state the agent may visit but will eventually leave
  \item[$Q$ matrix] The sub-matrix of $P$ containing only transient-to-transient transitions
  \item[Fundamental matrix] $N = (I-Q)^{-1}$; entry $N(i,j)$ = expected visits to state
    $j$ starting from state $i$
  \item[Expected steps] $\sum_j N(0,j)$: total expected tool calls before the run ends
  \item[KL divergence] A measure of information difference in bits; here, how much
    knowing the previous state changes the prediction of the next state
  \item[First-order Markov] Next state depends only on current state
  \item[Second-order Markov] Next state depends on current + previous state
  \item[JAR cluster] The JAR\_INSPECT self-loop; agent searching Maven cache for imports
  \item[FIX/BUILD loop] The rework cycle; agent editing tests and re-running
  \item[PRODUCTIVE flow] EXPLORE$\to$WRITE$\to$BUILD$\to$VERIFY; forward progress
  \item[TPE] Transition Probability Engineering; deliberately modifying $P$ to
    improve agent efficiency
  \item[$\Delta P$] Delta heatmap ($P_B - P_A$); shows which transitions changed
    between two variants
\end{description}

\end{document}
